% ============================================================
% 2. Related Systems
% ============================================================
\chapter{Related Systems}\label{ch:related}

Several existing applications address parts of the problem space covered by
CiboAmico. Table~\ref{tab:related} compares them along the dimensions that are
central to this project.

\begin{table}[H]
\centering
\caption{Comparison with related systems.}
\label{tab:related}
\begin{tabular}{@{}p{3.2cm}p{5.2cm}p{5.2cm}@{}}
\toprule
\textbf{System} & \textbf{Strengths} & \textbf{Limitations} \\
\midrule
Too Good To Go & Reduces food waste; large local network; simple ordering flow. & No pantry/inventory management; no recipe suggestions; no vendor publication control. \\
MyFoodRepo & Recipe database with ingredient filtering. & No local marketplace; no direct order from farmers; no substitution engine. \\
Supermarket apps (loyalty) & Shopping lists and discounts. & No integration with home inventory; no barcode-to-pantry flow. \\
\bottomrule
\end{tabular}
\end{table}

\section{Innovation Elements}

CiboAmico introduces three distinguishing elements with respect to the
systems above (design decisions D-01/D-02/D-03):

\begin{enumerate}
    \item \textbf{Ingredient substitution engine (D-01)}: when an ingredient is
    missing, the recipe engine proposes equivalent substitutes with a
    conversion factor (e.g.\ 0.5~kg of tomato purée for 1~kg of fresh
    tomatoes), implemented as a \emph{Strategy} pattern;
    \item \textbf{Dynamic role metamorfosi (D-02)}: a user can acquire the
    vendor or nutritionist role at runtime (Whole-Part pattern) without
    restructuring the class model;
    \item \textbf{Single direct order (D-03)}: the marketplace adopts a
    one-buyer--one-vendor direct order with no cart, tailored to the
    short-supply-chain scenario.
\end{enumerate}

\section{Design Rationale}

This section records the defence arguments for the three innovation
elements, as they are the most frequently questioned during the oral
examination.

\paragraph{D-01 (Ingredient substitution engine)} The substitution engine
implements the \emph{Strategy} pattern: each matching policy
(\texttt{StrictMatchingStrategy}, \texttt{SubstitutionMatchingStrategy}) is
an interchangeable algorithm, and new substitution policies (by price, by
calories, by proximity) can be added without modifying the existing code
(Open/Closed Principle). The engine is exposed through a single entry
point, the \texttt{RicettaMatchingFacade} (Facade pattern), which decouples
the presentation layer from the matching policies.

\paragraph{D-02 (Dynamic role metamorfosi)} The classical alternative to
role dynamism is rigid inheritance, which would force the system to
destroy the in-memory instance and create a new one (with different
references) every time a \texttt{Client} acquires the \texttt{Vendor} or
\texttt{Nutritionist} role. CiboAmico adopts the \emph{Whole-Part} pattern:
\texttt{Utente} (Whole) aggregates a collection of roles (Part), so the
role metamorphosis happens at runtime without altering the identity of
the object in memory. This avoids the Fragile Base Class problem and
keeps the model extensible.

\paragraph{D-03 (Single direct order)} The marketplace deliberately
avoids the cart idiom: the short-supply-chain scenario is one
buyer--one vendor per transaction. The order is created directly in the
\emph{CREATED} state and progresses through the state machine defined by
BR-04; the vendor is notified asynchronously through the
\emph{Observer} pattern (\texttt{VenditoreNotifier}).
